We treat graphs as simple undirected graphs (no loops, no multiple edges).
For a graph G=(V,E), the chromatic number \chi(G) is the least cardinal \kappa such that V can be written as a disjoint union of \kappa independent sets.

Ambiguity: the problem statement says “subgraph on \aleph_1 vertices”.
For an upper bound condition “\chi(\cdot)\le\aleph_0”, this is equivalent to requiring the induced subgraph on that vertex set to have \chi\le\aleph_0 (since deleting edges cannot increase chromatic number).
Accordingly, in what follows “subgraph on X\subseteq V” means the induced subgraph G[X].

Question (A): Does there exist a graph G with |V|=\aleph_2 and \chi(G)=\aleph_2 such that for every X\subseteq V with |X|=\aleph_1 we have \chi(G[X])\le \aleph_0?

Question (B): Does there exist a graph H with |V(H)|=\aleph_{\omega+1} and \chi(H)=\aleph_1 such that for every Y\subseteq V(H) with |Y|=\aleph_{\omega} we have \chi(H[Y])\le \aleph_0?

Edge cases/conventions: We use \aleph_0=|\mathbb N|. For induced subgraphs of size \aleph_1, “\chi\le \aleph_0” means “countably colorable”.
